\documentclass[class=article, crop=false, oneside]{standalone}
% ==============================================================================
% SETUP.TEX - CONFIGURACIÓN GLOBAL DEL PROYECTO TC2
% ==============================================================================
\ifdefined\SETUPCONFIGURED
\else
\newcommand{\SETUPCONFIGURED}{}

% Este archivo centraliza todos los paquetes y estilos.
% Debe incluirse en el preámbulo del libro principal y de los archivos standalone.

% --- 1. CODIFICACIÓN E IDIOMA ---
\usepackage[utf8]{inputenc}
\usepackage[T1]{fontenc}
\usepackage[spanish, es-tabla]{babel}  
\usepackage{csquotes} 

\usepackage{import}
\usepackage[mode=image, subpreambles=false]{standalone}

% --- 2. MATEMÁTICAS Y CIENCIAS ---
\usepackage{amsmath, amssymb, amsfonts, mathtools}

% Operadores en notación español/europea (seno, coseno y seno hiperbólicos)
\DeclareMathOperator{\sen}{sen}
\DeclareMathOperator{\ch}{ch}
\DeclareMathOperator{\sh}{sh}

% --- 3. DISEÑO Y GEOMETRÍA --- 
\usepackage[left=2.5cm,right=2.5cm,top=3cm,bottom=3cm]{geometry}
\usepackage{parskip} 
\usepackage{float}   

% --- 4. GRÁFICOS Y TABLAS ---
\usepackage{graphicx}
\usepackage{booktabs} 
\usepackage{tikz}
\usetikzlibrary{arrows.meta, calc, angles, quotes, babel, shapes.geometric}
\usepackage{pgfplots}
\usepgfplotslibrary{groupplots}
\usepgfplotslibrary{external}
\usepgfplotslibrary{patchplots}
\pgfplotsset{compat=1.18}
\usepackage[siunitx, RPvoltages]{circuitikz}

% --- 5. COLORES Y CAJAS ---
\usepackage[table]{xcolor}
\usepackage[theorems, most]{tcolorbox}

\definecolor{utnblue}{RGB}{0, 51, 102}
\definecolor{codegreen}{rgb}{0,0.6,0}
\definecolor{codegray}{rgb}{0.5,0.5,0.5}
\definecolor{codepurple}{rgb}{0.58,0,0.82}
\definecolor{backcolour}{rgb}{0.96,0.96,0.96}

% --- 6. ENTORNOS PERSONALIZADOS (CAJAS) ---
\newtcolorbox{pedagogico}{
    colback=codegreen!5!white,
    colframe=codegreen,
    title=\textbf{Enfoque Pedagógico},
    fonttitle=\bfseries,
    sharp corners,
    boxrule=0.5mm
}

\newtcolorbox{ingenieria}{
    colback=utnblue!5!white,
    colframe=utnblue,
    title=\textbf{Utilidad Ingenieril},
    fonttitle=\bfseries,
    sharp corners,
    boxrule=0.5mm
}

\newtcolorbox{nota}{
    colback=codegray!5!white,
    colframe=codegray,
    title=\textbf{Nota},
    fonttitle=\bfseries,
    sharp corners,
    boxrule=0.5mm
}

% --- 7. CÓDIGO FUENTE (LISTINGS) ---
\usepackage{listings}
\lstdefinestyle{miestilo}{
    backgroundcolor=\color{backcolour},   
    commentstyle=\color{codegreen},
    keywordstyle=\color{magenta},
    numberstyle=\tiny\color{codegray},
    stringstyle=\color{codepurple},
    basicstyle=\ttfamily\footnotesize,
    breakatwhitespace=false,         
    breaklines=true,                 
    captionpos=b,                    
    keepspaces=true,                 
    numbers=left,                    
    numbersep=5pt,                  
    showspaces=false,                
    showstringspaces=false,
    showtabs=false,                  
    tabsize=2,
    frame=single,                    
    rulecolor=\color{codegray}
}
\lstset{style=miestilo}
\lstset{literate={ñ}{{\~n}}1 {á}{{\'a}}1 {é}{{\'e}}1 {í}{{\'i}}1 {ó}{{\'o}}1 {ú}{{\'u}}1}

% --- 8. HIPERVÍNCULOS ---
\usepackage[colorlinks=true, linkcolor=utnblue, urlcolor=utnblue, citecolor=utnblue]{hyperref}

% --- 9. PLACEHOLDER PARA FIGURAS PENDIENTES ---
% Marca el lugar y la descripción de una figura aún no dibujada (p. ej. Cap. 5:
% circuitos y gráficas diferidos). Greppable por \figpendiente para el pase posterior.
\newcommand{\figpendiente}[1]{%
  \begin{center}%
  \setlength{\fboxsep}{6pt}%
  \fbox{\parbox{0.85\linewidth}{\centering\small\itshape
  [\textbf{Figura pendiente}]\\[2pt] #1}}%
  \end{center}}

\fi


% ------------------------------------------------------------------
% Trabajo Práctico Nº 4 — Filtros: Teoría Imagen
% PARTE I: Enunciado + respuestas de la guía.
% (La resolución completa está en TC-TP4-24-resolucion.tex.)
% Cátedra Teoría de Circuitos 2 — UTN-FRM — C. J. Larriqueta.
% ------------------------------------------------------------------

\begin{document}

% ===================== PORTADA / TÍTULO =====================
\begin{center}
    {\Large\bfseries Universidad Tecnológica Nacional --- F.R. Mendoza}\\[2pt]
    {\large Teoría de Circuitos 2}\\[10pt]
    {\huge\bfseries\color{utnblue} Trabajo Práctico Nº 4}\\[4pt]
    {\LARGE\bfseries Filtros: Teoría Imagen}\\[10pt]
    {\large Enunciado y respuestas}\\[2pt]
\end{center}

\vspace{0.5ex}
\hrule
\vspace{1ex}

% ============================================================
% ============== ENUNCIADO ====================================
% ============================================================
\section*{\color{utnblue}Enunciado}

\begin{enumerate}
    \item Esquematizar el diseño de un filtro completo según la Teoría Imagen.

    \item Calcular el prototipo de un filtro pasabajo de configuración ``T'' con
    los siguientes parámetros:
    \begin{equation*}
        \left\{
        \begin{aligned}
            \omega_c &= 20.000 \ \tfrac{\text{rad}}{\text{s}}\\
            R_L &= 500 \ \Omega
        \end{aligned}
        \right.
    \end{equation*}

    \item Calcular el ejercicio anterior partiendo del circuito normalizado.

    \item Tomando el prototipo del ejercicio Nº 2, diseñar la sección
    $m$-derivada para una frecuencia de atenuación infinita
    $\omega_\infty = 22.500 \ \tfrac{\text{rad}}{\text{s}}$.

    \item Calcular la hemi-sección de adaptación correspondiente al ejercicio
    anterior.

    \item A partir del filtro pasabajos obtener, mediante la transformación en
    frecuencia:
    \begin{enumerate}
        \item[a)] un pasaaltos, especificando la transformación de elementos, la
        frecuencia de corte, el ancho de banda y la impedancia característica;
        \item[b)] un pasabanda.
    \end{enumerate}

    \item El circuito de la figura representa el prototipo de un filtro pasa
    altos imagen, donde $C = 10^{-7}\ \text{F}$:
    \begin{figure}[H]
        \centering
        % Figura generada con netlist2tikz (n2t render fig_ej07_enunciado.sch --tikz --no-standalone).
% Netlist fuente (fig_ej07_enunciado.sch):
%   W in_t 1; right=0.4
%   C1 1 2; right=2, l^=C
%   C2 2 3; right=2, l^=C
%   W 3 out_t; right=0.4
%   L 2 0; down=1.6, l_=L
%   W in_b 0; right=2.4
%   W 0 out_b; right=2.4
%
\begin{tikzpicture}[scale=1.00, transform shape, /tikz/circuitikz/bipoles/length=1.50cm, american currents, american voltages, font={\fontsize{7.5}{9}\selectfont}, voltage dir=RP]
  \coordinate (in_t) at (0,1.92);
  \coordinate (1) at (0.48,1.92);
  \coordinate (2) at (2.88,1.92);
  \coordinate (3) at (5.28,1.92);
  \coordinate (out_t) at (5.76,1.92);
  \coordinate (0) at (2.88,0);
  \coordinate (in_b) at (0,0);
  \coordinate (out_b) at (5.76,0);
  \draw (in_t) to (1);
  \draw (1) to [C, l^=C, n=C1] (2);
  \draw (2) to [C, l^=C, n=C2] (3);
  \draw (3) to (out_t);
  \draw (2) to [L, l_=L, n=L] (0);
  \draw (in_b) to (0);
  \draw (0) to (out_b);
  \draw (in_t) node[ocirc] {};
  \draw (out_t) node[ocirc] {};
  \draw (in_b) node[ocirc] {};
  \draw (out_b) node[ocirc] {};
\end{tikzpicture}

    \end{figure}
    Calcular:
    \begin{enumerate}
        \item[a)] el valor del elemento ``$L$'' sabiendo que la impedancia
        característica $z_{0T}(\infty) = 100\ \Omega$ cuando $\omega\to\infty$;
        \item[b)] la frecuencia de corte $\omega_c$.
    \end{enumerate}

    \item Dado el circuito $m$-derivada normalizado de un filtro pasaalto
    (designados como $C_1'$, $L'$, $C_2'$), calcular la hemisección. Determinar
    los valores normalizados.
    \begin{figure}[H]
        \centering
        % Figura generada con netlist2tikz (n2t render fig_ej08_enunciado.sch --tikz --no-standalone).
% Netlist fuente (fig_ej08_enunciado.sch):
%   W in_t 1; right=0.4
%   C1 1 2; right=2, l^=C_1'
%   C2 2 3; right=2, l^=C_1'
%   W 3 out_t; right=0.4
%   L 2 4; down=1.1, l_=L'
%   C3 4 0; down=1.1, l_=C_2'
%   W in_b 0; right=2.4
%   W 0 out_b; right=2.4
%
\begin{tikzpicture}[scale=1.00, transform shape, /tikz/circuitikz/bipoles/length=1.50cm, american currents, american voltages, font={\fontsize{7.5}{9}\selectfont}, voltage dir=RP]
  \coordinate (in_t) at (0,2.64);
  \coordinate (1) at (0.48,2.64);
  \coordinate (2) at (2.88,2.64);
  \coordinate (3) at (5.28,2.64);
  \coordinate (out_t) at (5.76,2.64);
  \coordinate (4) at (2.88,1.32);
  \coordinate (0) at (2.88,0);
  \coordinate (in_b) at (0,0);
  \coordinate (out_b) at (5.76,0);
  \draw (in_t) to (1);
  \draw (1) to [C, l^=$C_1'$, n=C1] (2);
  \draw (2) to [C, l^=$C_1'$, n=C2] (3);
  \draw (3) to (out_t);
  \draw (2) to [L, l_=L', n=L] (4);
  \draw (4) to [C, l_=$C_2'$, n=C3] (0);
  \draw (in_b) to (0);
  \draw (0) to (out_b);
  \draw (in_t) node[ocirc] {};
  \draw (out_t) node[ocirc] {};
  \draw (in_b) node[ocirc] {};
  \draw (out_b) node[ocirc] {};
\end{tikzpicture}

    \end{figure}
    \begin{equation*}
        C_1' = \frac{2}{m L_1}\ ;\qquad
        C_2' = \frac{4m}{L_1(1-m^2)}\ ;\qquad
        L' = \frac{1}{m C_2}
    \end{equation*}
    Tener en cuenta que los valores del pasabajos normalizado son:
    $L_1 = 2$; $C_2 = 2$ y $\omega_\infty = \sqrt{1-m^2}$.

    \item Calcular un filtro completo pasabanda con $R_L = 600\ \Omega$;
    $\omega_1 = 45.000\ \tfrac{\text{rad}}{\text{s}}$;
    $\omega_2 = 90.000\ \tfrac{\text{rad}}{\text{s}}$;
    $\omega_{\infty 1} = 38.000\ \tfrac{\text{rad}}{\text{s}}$.

    \item Calcular un filtro completo suprimebanda con $R_L = 75\ \Omega$;
    $f_1 = 1.000\ \text{Hz}$; $f_2 = 3.500\ \text{Hz}$; $f_{\infty 2} = 3.100\ \text{Hz}$.
\end{enumerate}

% -------------------- RESPUESTAS OFICIALES --------------------
\subsection*{Respuestas (guía)}

\begin{description}
    \item[1)] Esquema de filtro completo (ver resolución).
    \item[2)] $L_1 = 50\ \text{mH}$; $C_2 = 200\ \text{nF}$. Los valores finales
    de la rama serie se obtienen dividiendo $L_1$ por dos.
    \item[4)] $L_1 = 23\ \text{mH}$; $L_2 = 21{,}6\ \text{mH}$;
    $C_2 = 91{,}6\ \text{nF}$. Ídem anterior.
    \item[5)] $L_1 = 30\ \text{mH}$; $L_2 = 26{,}7\ \text{mH}$;
    $C_2 = 60\ \text{nF}$. Ídem anterior.
    \item[6a)] Transformación $s\to\dfrac{1}{s'}$; elementos
    $C_1 = \dfrac{2}{L_1}$, $L_2 = \dfrac{1}{C_2}$; frecuencia de corte
    $\omega_c = \dfrac{1}{2\sqrt{L_2 C_1}}$; impedancia característica en
    $\omega\to\infty$: $z_{0T(\infty)} = \sqrt{\dfrac{L_2}{C_1}}$; ancho de
    banda $AB:(\omega_c\ ;\ \infty)$.
    \item[7)] a) $L = 500\ \mu\text{H}$; \qquad b) $\omega_c = 100.000\ \tfrac{\text{rad}}{\text{s}}$.
    \item[8)] Hemisección de entrada del pasaaltos:
    $C_s = C_1' = \dfrac{2}{m L_1} = 1{,}67\ \text{F}$;
    $L_d = 2L' = \dfrac{2}{m C_2} = 1{,}67\ \text{H}$;
    $C_d = \dfrac{C_2'}{2} = \dfrac{2m}{(1-m^2)L_1} = 0{,}94\ \text{F}$.
    \item[9)] Prototipo: $L_1 = 13\ \text{mH}$; $C_1 = 18{,}5\ \text{nF}$;
    $L_2 = 3{,}3\ \text{mH}$; $C_2 = 74\ \text{nF}$.\\
    $m$-Derivada: $L_1 = 10\ \text{mH}$; $C_1 = 24{,}5\ \text{nF}$;
    $L_2 = 4{,}4\ \text{mH}$; $C_2 = 56\ \text{nF}$; $L_3 = 3{,}8\ \text{mH}$;
    $C_3 = 65\ \text{nF}$.\\
    Hemisección de adaptación: $L_1 = 8\ \text{mH}$; $C_1 = 31\ \text{nF}$;
    $L_2 = 11{,}1\ \text{mH}$; $C_2 = 22{,}2\ \text{nF}$; $L_3 = 14{,}2\ \text{mH}$;
    $C_3 = 17{,}4\ \text{nF}$.
    \item[10)] Prototipo: $L_1 = 8{,}5\ \text{mH}$; $C_1 = 849\ \text{nF}$;
    $L_2 = 2{,}4\ \text{mH}$; $C_2 = 3.030\ \text{nF}$.\\
    Sección $m$-Derivada: $L_1 = 5{,}3\ \text{mH}$; $C_1 = 1.380\ \text{nF}$;
    $L_2 = 3{,}9\ \text{mH}$; $C_2 = 1.860\ \text{nF}$; $L_3 = 4{,}3\ \text{mH}$;
    $C_3 = 1.680\ \text{nF}$.\\
    Hemisección de adaptación: $L_1 = 5{,}1\ \text{mH}$; $C_1 = 1.415\ \text{nF}$;
    $L_2 = 9{,}1\ \text{mH}$; $C_2 = 796\ \text{nF}$; $L_3 = 8{,}0\ \text{mH}$;
    $C_3 = 910\ \text{nF}$.
\end{description}

\end{document}
